% Image credits for the photographs and AI illustrations of Book 5
% (University Biology, Year 3). Machine-readable license records live in
% images/book5/CREDITS.md; keep the two files in sync.
\chapter*{Crédits des images}

% \sloppy must be in force when the paragraph is BROKEN, hence the \par before
% \endgroup -- without it the group closes first and the tolerance is restored.
\begingroup\sloppy

Les dessins de ce livre sont originaux. Les photographies sont
utilisées sous leurs licences libres respectives, avec nos
remerciements à leurs auteurs~:

\begin{itemize}
  \item Chapitre~\ref{ch:b3:genomics}~: photographie de Frederick
        Sanger --- domaine public (Wikimedia Commons).
  \item Chapitre~\ref{ch:b3:genetic-engineering}~: riz doré à côté de
        riz ordinaire, International Rice Research Institute, 2011 ---
        CC~BY~2.0 (Wikimedia Commons).
  \item Chapitre~\ref{ch:b3:structural-biology}~: photographie de Max
        Perutz, Associated Press, 1962 --- domaine public (Wikimedia
        Commons).
  \item Chapitre~\ref{ch:b3:bacteriology}~: portrait de Robert Koch,
        Wellcome Collection --- CC~BY~4.0 (Wikimedia Commons)~;
        photographie d'Alexander Fleming dans son laboratoire, archives
        du U.S. Navy Bureau of Medicine and Surgery --- aucune
        restriction de droits connue (Wikimedia Commons).
  \item Chapitre~\ref{ch:b3:innate-immunity}~: photographie d'Élie
        Metchnikoff, Bain News Service, Library of Congress --- domaine
        public (Wikimedia Commons).
  \item Chapitre~\ref{ch:b3:adaptive-immunity}~: portrait gravé
        d'Edward Jenner, 1807 --- domaine public (Wikimedia Commons).
  \item Chapitre~\ref{ch:b3:nervous-systems}~: dessin d'une cellule de
        Purkinje par Santiago Ramón y Cajal, 1899 --- domaine public
        (Wikimedia Commons).
  \item Chapitre~\ref{ch:b3:endocrinology}~: photographie de Frederick
        Banting et Charles Best, vers 1924 --- domaine public (Wikimedia
        Commons).
  \item Chapitre~\ref{ch:b3:molecular-evolution}~: photographie de Motoo
        Kimura, 1986, photographe inconnu, via S.~Kumar et T.~Gojobori,
        \emph{Molecular Biology and Evolution} (2023) --- CC~BY~4.0
        (Wikimedia Commons).
  \item Chapitre~\ref{ch:b3:behavioural-ecology}~: photographie de Niko
        Tinbergen par Rob Mieremet, Anefo, 1973 --- CC0, Nationaal
        Archief (Wikimedia Commons).
  \item Chapitre~\ref{ch:b3:conservation-biology}~: photographie de
        Martha, la dernière tourte voyageuse, Enno Meyer, 1912 ---
        domaine public (Wikimedia Commons)~; loups du parc national de
        Yellowstone, National Park Service, 1996 --- domaine public
        (Wikimedia Commons)~; Sirocco le kakapo, Department of
        Conservation, Nouvelle-Zélande, 2012 --- CC~BY~2.0 (Wikimedia
        Commons).
\end{itemize}

Les liens vers les sources complètes et les adresses des licences de
chaque photographie sont donnés dans
\texttt{images/\allowbreak book5/\allowbreak CREDITS.md}, dans le dépôt du livre.

\bigskip
À côté des photographies et des dessins originaux, quelque
quatre-vingt-dix figures réparties dans les chapitres du volume
sont des illustrations produites par intelligence artificielle,
réalisées avec la génération d'images d'OpenAI à partir d'invites
écrites par les éditeurs du livre, et relues pour en vérifier
l'exactitude biologique avant leur inclusion. L'invite complète de chaque image générée est
consignée dans \texttt{images/\allowbreak book5/\allowbreak ai/\allowbreak PROMPTS.md},
dans le dépôt.
\par\endgroup
\clearpage
